\section*{第 \ref{ch:g4:geometry} 章——直线与多边形}

\begin{solution}{exo:g4:geometry:1}
作图用到\cref{met:g4:geometry:draw}；小方框
表示 $d$ 与垂线之间的\omterm{def:g3:shapes:rightangle}{直角}。
\end{solution}

\begin{solution}{exo:g4:geometry:2}
E：三条水平线段互相\omterm{def:g4:geometry:def}{平行}；每一条都与竖线\omterm{def:g4:geometry:def}{垂直}。
H：两条竖线段互相\omterm{def:g4:geometry:def}{平行}，
横线与两条竖线都\omterm{def:g4:geometry:def}{垂直}。N：两条竖线段互相\omterm{def:g4:geometry:def}{平行}；
斜线两者都不是。T：两条线段互相\omterm{def:g4:geometry:def}{垂直}。
Z：上下两条线段互相\omterm{def:g4:geometry:def}{平行}；
斜线两者都不是。
\end{solution}

\begin{solution}{exo:g4:geometry:3}
\omterm{def:g4:geometry:polygon}{六边形}；\omterm{def:g4:geometry:polygon}{四边形}；三角形；\omterm{def:g4:geometry:polygon}{八边形}；\omterm{def:g4:geometry:polygon}{五边形}。\omterm{def:g4:geometry:polygon}{多边形}的
\omterm{def:g4:geometry:polygon}{顶点}个数与边数相等：$6$、$4$、$3$、$8$、$5$。
\end{solution}

\begin{solution}{exo:g4:geometry:4}
\omterm{def:g4:geometry:polygon}{五边形}有 $5$ 条边。从 $A$ 出发，不相邻的\omterm{def:g4:geometry:polygon}{顶点}是 $C$ 和
$D$：两条对角线，$[AC]$ 与 $[AD]$。
\end{solution}

\begin{solution}{exo:g4:geometry:5}
$OA = 4$ cm（\omterm{def:g3:shapes:shapes}{半径}）；$OM = 4$ cm（圆上每一点到圆心的距离）；
$AB = 8$ cm（\omterm{def:g4:geometry:circle}{直径}是\omterm{def:g3:shapes:shapes}{半径}的两倍）。
\end{solution}

\begin{solution}{exo:g4:geometry:6}
先沿一条边张开圆规，再把这个张口移到另外两条边上：
若圆规两脚的距离恰好与第二条边吻合，
则这两条边相等。
\end{solution}

\begin{solution}{exo:g4:geometry:7}
这条路径围成一个有 $5$ 条边的\omterm{def:g4:geometry:polygon}{五边形}——
但不是正\omterm{def:g4:geometry:polygon}{五边形}！
\end{solution}

\begin{solution}{exo:g4:geometry:8}
“\omterm{def:g4:geometry:def}{垂直}的两条直线一定相交”：正确——
它们恰好在\omterm{def:g3:shapes:rightangle}{直角}处相交。

“\omterm{def:g4:geometry:def}{平行}的两条直线永不相交”：正确——这正是\omterm{def:g4:geometry:def}{平行}的定义。

“一条直线可以既与某条直线\omterm{def:g4:geometry:def}{平行}，又与另一条直线\omterm{def:g4:geometry:def}{垂直}”：
正确——在方格纸上，一条水平线与每一条水平线都\omterm{def:g4:geometry:def}{平行}，
并与每一条竖线都\omterm{def:g4:geometry:def}{垂直}。
\end{solution}

\begin{solution}{exo:g4:geometry:9}
直线 $d$ 与 $f$ 看上去\omterm{def:g4:geometry:def}{平行}——而且它们确实总是\omterm{def:g4:geometry:def}{平行}的：
\omterm{def:g4:geometry:def}{垂直}于同一条直线的两条直线互相\omterm{def:g4:geometry:def}{平行}
（由\cref{prop:g6:lines:facts}正式确认这一点）。
\end{solution}

\begin{solution}{exo:g4:geometry:10}
一种正确的作图步骤是：“1. 画一个正方形 $ABCD$，边长为 $4$ cm。
2. 标出两条对角线的交点 $O$。3. 画圆，
圆心为 $O$、\omterm{def:g3:shapes:shapes}{半径}为 $2$ cm。”这个圆与每条边都相切，
切点恰为各边的中点。
\end{solution}

\begin{solution}{exo:g4:geometry:11}
\omterm{def:g4:geometry:polygon}{四边形}：$2$ 条对角线。\omterm{def:g4:geometry:polygon}{五边形}：$5$ 条。\omterm{def:g4:geometry:polygon}{六边形}：$9$ 条。每增加一个\omterm{def:g4:geometry:polygon}{顶点}，
对角线的条数增加得越来越多——增长越来越快
（$+3$，然后 $+4$，\dots）。
\end{solution}
